% =====================================================================
% Chapter 2: Literature Review  --  DRAFT v1
% Project: Vision-Language Models for Automated Testing of Dynamic
%          Interactive UI and Animation Anomalies
%
% ---------------------------------------------------------------------
% 【使用說明 / HOW TO USE】
%
% 標註等級 (annotation levels):
%   % [OK]      = Claude 已親自讀過原文並確認，數字/主張可信
%   % [VERIFY]  = 引用資訊來自二手來源（他人參考書目 / 你的專案筆記），
%                 書目細節大致可信但你必須點開原文確認
%   % [!! RED]  = Claude 完全沒有查證過，可能根本不存在或名稱有誤，
%                 在你親自確認之前，絕對不要送出
%
% 中文註解若造成 pdflatex 編譯問題（極少見），可用下列指令一次刪除：
%   sed -i '' '/^% \[/d' chapter2_literature_review.tex
% =====================================================================

\chapter{Literature Review}
\label{ch:litreview}

% ---------------------------------------------------------------------
% 2.1 Chapter roadmap
% 依「論文書寫語料」要求：章節開頭必須有明確的引導語
% ---------------------------------------------------------------------
\section{Overview}
\label{sec:lit-overview}

This chapter reviews the body of work on which the present study builds.
Before turning to the use of vision-language models (VLMs) for user
interface (UI) defect detection, it is important first to introduce some
of the key concepts relating to visual regression testing (VRT) and the
oracle problem it inherits, since the limitations of the conventional
pixel-level oracle constitute the practical motivation for this work.
Section~\ref{sec:vrt} therefore examines VRT and the false-positive
problem associated with pixel-level comparison. Section~\ref{sec:temporal}
then considers how the field has characterised web UI visual defects, and
argues that the temporal class of defects remains inadequately covered.
Section~\ref{sec:vlm-se} surveys the emerging application of VLMs to
visual software engineering tasks and identifies a recurring failure mode.
Section~\ref{sec:input} reviews how video input is constructed for such
models, and Section~\ref{sec:datasets} discusses the datasets and
synthetic fault injection techniques available for evaluation.
Section~\ref{sec:metrics} addresses evaluation criteria. Finally,
Section~\ref{sec:gap} synthesises these strands into an explicit statement
of the research gap that this dissertation addresses.


% ---------------------------------------------------------------------
% 2.2 VRT and the pixel oracle
% ---------------------------------------------------------------------
\section{Visual Regression Testing and the Pixel-Level Oracle}
\label{sec:vrt}

\subsection{The pixel-diff oracle}
\label{sec:vrt-oracle}

Visual regression testing is a standard quality assurance step in modern
software release pipelines. On each change, UI screenshots are re-rendered
and compared against an approved baseline image, and any detected
difference is routed to a human reviewer who decides whether it represents
an intended update or an unintended regression~\citep{zhang2026beyond}.
% [OK] Zhang et al. 2026 — Claude 已讀過全文，此段落敘述與原文一致。

The dominant comparison strategy, particularly in open-source and
continuous-integration pipelines, is pixel-level comparison. Its appeal is
that it requires no specification of expected appearance and no
instrumentation of the application under test; its weakness is that it is
semantically blind. It flags any deviation, treating sub-pixel rendering
jitter, anti-aliasing differences and dynamic content updates identically
to genuine structural regressions such as missing components or layout
restructuring~\citep{zhang2026beyond}.

The consequence is a false-positive rate high enough to threaten the
practical viability of the technique. \citet{zhang2026beyond}
quantify this directly on WUICC-bench, a benchmark of 9{,}906
human-verified Web UI change samples of which 1{,}323 represent
non-meaningful visual noise. A naive pixel diff correctly reports
``no change'' on \SI{0}{\percent} of the non-meaningful samples,
corresponding to a false-positive rate of \SI{100}{\percent}; adding a
0.1 difference threshold with anti-aliasing tolerance raises suppression
only to \SI{6.72}{\percent}, a false-positive rate of
\SI{93.28}{\percent}, while already beginning to miss genuine changes.
By contrast, trained image-difference-captioning models on the same
benchmark suppress between \SI{85.82}{\percent} and \SI{97.39}{\percent}
of non-meaningful changes while still reporting over \SI{93}{\percent} of
genuine ones. These figures establish that the deficiency of the pixel
oracle is not a matter of parameter tuning but of the representation
itself: a single fixed threshold cannot separate noise from signal.
% [OK] 上述所有數字皆出自 Zhang et al. 2026, Table VI, Claude 已核對。
% [VERIFY] 但請注意兩點必須在正文或註腳交代：
%   (a) 該文為 arXiv preprint (v1, 2026-07-02)，尚未經同儕審查；
%   (b) 此處 FPR 定義為「在人工注入之非實質變更上未回報 no change 的比例」，
%       與產線 CI 中「工程師實際分流掉的告警比例」並非同一操作型定義。
%       審查員很可能會問這一題，你必須能區分。

\subsection{Attempts to suppress false positives}
\label{sec:vrt-suppression}

Much of the research in this area has focused on reducing the manual
confirmation burden that pixel comparison imposes rather than on replacing
the oracle. \citet{adachi2018reducing} group identical
differences so that a reviewer confirms each distinct difference only once.
\citet{tanno2020region} propose region-based matching to
distinguish essential from inessential differences. A subsequent method by
\citet{adachi2020masking} masks dynamic content areas by
exploiting the positional relationships between screen elements, and
MIXVRT~\citep{amano2021mixvrt} combines image comparison with HTML code
comparison.
% [VERIFY] 上述四筆的作者/年份/會議名稱，我是從 Zhang et al. 2026 的
%   參考書目轉抄的，沒有親自讀過任何一篇。你必須：
%   1. 上 IEEE Xplore / J-STAGE 確認書目資訊正確
%   2. 確認我對每一篇「做了什麼」的一句話描述沒有寫錯
%   3. adachi2018reducing 是 2 頁 poster，內容可能比我描述的更薄，
%      引用時語氣要保守
% [!! RED] "amano2021mixvrt" 這個 key 的第一作者我是猜的 —— Zhang et al.
%   的參考書目編號 [7] 我沒有取到完整作者名。請務必查出真正作者再引用，
%   或直接刪掉這一句。

Two observations follow. First, all of these techniques output a binary
flag or a highlighted region, and therefore reduce rather than eliminate
the reviewer's interpretive burden. Second, and more importantly for the
present study, all of them operate on pairs of static screenshots. None
addresses defects whose evidence exists only across time.


% ---------------------------------------------------------------------
% 2.3 Temporal defects
% ---------------------------------------------------------------------
\section{From Static Screenshots to Temporal UI Defects}
\label{sec:temporal}

\subsection{Taxonomies of web UI visual defects}
\label{sec:temporal-taxonomy}

Having discussed the mechanics of VRT, the following section considers
what is actually being detected. There is some agreement in the literature
that web UI visual defects fall into a small number of recurring
categories. \citet{zhang2026beyond} operationalise a 37-rule
taxonomy comprising nine meaningful categories --- missing elements,
added elements, attribute modification, layout changes, reordering,
resizing, content update, thematic changes and replacement --- and three
non-meaningful categories representing visual noise, namely geometric
shifts, dynamic content changes and pure style changes.
% [OK] 此分類法我已從原文 Table I 核對。

A complementary industrial taxonomy distinguishes layout regression,
colour or style change, content truncation, interactive state anomalies,
typography errors and responsive layout failures~\citep{vrtatscale}.
Behavioural anomaly taxonomies, which characterise interruptions to a task
flow rather than to a rendered page, add categories such as action
failure, unexpected authentication interruption, obstructive pop-ups,
loading delay and network disconnection~\citep{guirobust}.
% [!! RED] 這兩筆是本章最高風險的引用。
%   - "Visual Regression Testing at Scale"：這個標題只出現在你的第二次
%     會議筆記裡，我在網路上完全沒有查證過。作者、年份、發表場域全部未知，
%     連這篇論文是否存在、標題是否正確我都不能確定。
%   - "GUIRobust (KDD 2026)"：同樣只出現在你的筆記裡。你的筆記還提到它
%     宣稱把標註成本降低 19 倍 —— 這個數字我也沒查證。
%   → 這兩篇是你自己找到的，請你直接把 PDF 或 DOI 找出來核對書目，
%     並確認我上面複述的「六類」與「七類」內容沒有被筆記轉述時扭曲。
%     若查不到原文，寧可整段刪掉，改成「本研究自行定義之缺陷分類」。

\subsection{Why temporal defects escape existing oracles}
\label{sec:temporal-gap}

The taxonomies above share a common assumption: that a defect is
observable in a single rendered state. Yet an important class of modern
web UI defects is not. Cumulative layout shift during asynchronous
loading, component flickering caused by redundant re-rendering,
transition stuttering and animations that freeze mid-flight all manifest
as a discontinuity between states rather than as a property of any one
state. Two consequences follow for automated testing. First, a
screenshot-based oracle sampled at an arbitrary instant may capture the
interface either before or after the anomalous moment and report nothing
wrong. Second, DOM-based assertions in tools such as Playwright or
Selenium may pass, because the structural state eventually converges on
the correct value even when the visual trajectory towards it is defective.
It is important that studies into this class of defect be carried out
because web applications, unlike the entertainment software in which
temporal glitch detection has principally been studied, incur direct
commercial cost when interaction is disrupted.
% [VERIFY] 這一小節是論證，不是文獻回顧，因此沒有引用。這是可以的，
%   但審查員可能會要求你替「DOM 斷言會通過但視覺仍然壞掉」這句話補一個
%   引用或一個具體例子。建議你用自己資料集裡的 displaynone_collapse
%   或 skew_distortion 當例子，並在此處前向引用 Chapter 4。


% ---------------------------------------------------------------------
% 2.4 VLMs
% ---------------------------------------------------------------------
\section{Vision-Language Models for Visual Software Engineering}
\label{sec:vlm-se}

\subsection{Static visual tasks}
\label{sec:vlm-static}

One branch of the literature has begun applying VLMs to visual software
engineering tasks. Recent work has targeted the detection of cross-browser
visual inconsistencies~\citep{xbidetective} and the repair of visual
software issues through cross-modal reasoning~\citep{seeingisfixing}.
\citet{zhang2026beyond} extend this direction by reframing
regression checking as a change captioning problem, so that the tester is
told what changed in natural language rather than being handed a
highlighted region to interpret.
% [VERIFY] xbidetective 與 seeingisfixing 兩筆同樣是從 Zhang et al. 的
%   參考書目轉抄，我沒讀過原文，且沒有取到完整的作者與年份。
%   請查出完整書目，或若與你的論述關聯不強，直接刪除。

\subsection{The sensitivity--conservatism failure mode}
\label{sec:vlm-failure}

A consistent finding across this literature is that off-the-shelf VLMs do
not resolve the false-positive problem by themselves; they relocate it.
\citet{zhang2026beyond} evaluate two open general-purpose VLMs
zero-shot on WUICC-bench and report sharply asymmetric failures.
Qwen2-VL-7B-Instruct suppresses almost all non-meaningful changes,
reaching \SI{99.63}{\percent} no-change accuracy, but achieves this by
under-reporting genuine ones, with change accuracy falling to
\SI{65.68}{\percent}. Llama-3.2-11B-Vision-Instruct fails in the opposite
direction, with no-change accuracy of only \SI{21.64}{\percent}. Both
change accuracies fall well below those of the trained methods,
indicating that off-the-shelf VLMs do not yet provide the selective
suppression that VRT requires.
% [OK] 數字出自 Zhang et al. 2026 Table VI，已核對。

This pattern --- a model that is either excessively sensitive or
excessively conservative, but not appropriately calibrated --- has also
been reported in the temporal setting. TempGlitch, a benchmark for
temporal defect detection in gameplay video, concludes that VLMs supplied
with raw video perform close to chance and fall into precisely these two
extremes, and that denser uniform sampling does not remedy the
problem~\citep{tempglitch}.
% [!! RED] ★★★ 本章最關鍵、也最危險的一筆引用 ★★★
%   "TempGlitch" 是你整篇論文的出發點，但你給我的所有專案檔案裡，
%   從頭到尾都只有這個名字 —— 沒有作者、沒有年份、沒有發表場域、沒有 DOI。
%   我沒有查證過它，也不確定這是正式論文標題還是你筆記裡的簡稱。
%   你必須做的事（今天就做）：
%     1. 找出原始 PDF，確認正式標題、作者、年份、venue
%     2. 親自確認「near chance」「overly sensitive / overly conservative」
%        「denser sampling 無效」這三個結論是否真的出自該文，
%        以及它是否真的只做 gameplay
%     3. 確認它是不是 benchmark-only（你的筆記說是，但要自己看過）
%   → 這一筆若在 Demo 面談被問到而你答不出來，殺傷力極大，
%     因為你的整個 research gap 都架在它上面。


% ---------------------------------------------------------------------
% 2.5 Input construction
% ---------------------------------------------------------------------
\section{Input Construction for Video-Language Models}
\label{sec:input}

\subsection{Uniform sampling and its limitations}
\label{sec:input-uniform}

Because current VLMs accept a bounded number of image tokens, video must
be reduced to a frame sequence before it can be analysed. The simplest
strategy, and the one adopted by TempGlitch, is uniform temporal sampling
at a fixed rate~\citep{tempglitch}.
% [!! RED] 同上，TempGlitch 的取樣策略（你的筆記說是 FFmpeg 1 FPS / 5 FPS）
%   我完全沒有查證。請確認後再寫死具體幀率。

Uniform sampling is attractive for its simplicity but is poorly matched to
the statistics of UI interaction video, in which long static intervals are
punctuated by brief transitions. The majority of uniformly sampled frames
are therefore redundant, consuming tokens without contributing evidence,
while the short interval in which a temporal defect is observable may be
missed entirely.

\subsection{Keyframe selection}
\label{sec:input-keyframe}

An alternative is to select frames adaptively according to a measure of
inter-frame change. Classical approaches compute frame differencing or
optical flow to identify moments of significant motion. Applying such
measures to UI video is attractive because rendering is deterministic and
the background is static, so a change signal is comparatively
unambiguous compared with camera-captured footage.
% [VERIFY] 這一段目前完全沒有引用，太薄。你需要補 2--3 筆關鍵幀擷取 /
%   video summarisation 的經典文獻（例如 shot boundary detection 或
%   長影片 VLM 的 frame selection 研究）。這是我完全沒有替你查過的區塊。

\subsection{Structural signals as a complementary modality}
\label{sec:input-dom}

A property peculiar to the web domain is that a second, non-visual
description of the same interface is available at no cost: the DOM.
Hybrid approaches that combine image comparison with HTML comparison have
been explored in the VRT setting~\citep{amano2021mixvrt}, but principally
as a means of localising differences rather than as evidence supplied to a
reasoning model. The present study treats structural mutation events as a
first-class signal, used both to select frames and to constrain the final
decision.
% [VERIFY] 見前述 amano2021mixvrt 的紅旗。若查不到，改寫為
%   「hybrid image--HTML comparison has been explored」並引用 Zhang et al.
%   的 related work 二手轉述（但這是次佳做法，審查員可能會扣分）。


% ---------------------------------------------------------------------
% 2.6 Datasets
% ---------------------------------------------------------------------
\section{Datasets and Synthetic Fault Injection}
\label{sec:datasets}

Evaluating a UI anomaly detector requires paired normal and defective
interaction recordings with reliable ground truth, and such data is
scarce. GUI-World provides a large corpus of real GUI interaction videos
with annotations, but is oriented towards agent capability evaluation
rather than defect detection, and consequently contains few labelled
defective samples~\citep{guiworld}.
% [VERIFY] GUI-World 我沒有讀過原文，只有你筆記裡的 URL
%   (https://gui-world.github.io/) 與 2,499 支影片的數字。
%   請查出正式書目（作者、年份、venue、arXiv ID）並確認
%   「2,499 支正常網頁互動影片」這個數字的正確定義。

Manual production of defective recordings is prohibitively
labour-intensive, and this has motivated synthetic fault injection, in
which defects are introduced programmatically into a working application.
Recent work has demonstrated that automated pipelines can reduce
annotation cost by a substantial factor relative to manual
labelling~\citep{guirobust}, and \citet{zhang2026beyond}
construct their benchmark by an LLM-driven mutation pipeline that applies
exactly one atomic change per sample, retaining \SI{78.09}{\percent} of
generated samples after human verification with inter-annotator agreement
of $\kappa \approx 0.722$.
% [OK] 78.09% 與 kappa=0.722 出自 Zhang et al. 2026，已核對。
% [!! RED] 我刻意把你筆記裡的「19 倍」寫成模糊的 "a substantial factor"，
%   因為 GUIRobust 與那個數字我都沒查證。你查證後再決定要不要寫死數字。

The approach adopted in the present study differs from both. Rather than
mutating static HTML, defects are injected into a live application during
recorded interaction, so that the resulting anomaly is temporal in nature.
% [VERIFY] 請確認這句話準確描述你 Playwright 腳本實際做的事。


% ---------------------------------------------------------------------
% 2.7 Metrics
% ---------------------------------------------------------------------
\section{Evaluation Criteria}
\label{sec:metrics}

There is some agreement in the literature that overall accuracy is an
inadequate summary for detection tasks of this kind, and that the false
alarm rate on the negative class is the metric that governs practical
adoption. This position is well established in vision-based anomaly
detection more broadly: \citet{tao2023psnet}, for instance,
identify high false alarm rates as one of the two principal obstacles to
deployment in their domain, and report false alarm rate alongside
detection rate and F1-score rather than accuracy alone.
% [VERIFY] Tao et al. 2023 (Pattern Recognition, 森林煙霧辨識) 我已讀過，
%   書目與該主張皆正確。但這是跨領域類比，與 Web UI 沒有直接關係。
%   ※ 這一段是「可選」的。如果你覺得跨領域引用會讓審查員覺得牽強，
%     整段刪掉不影響論證。若保留，務必寫成「跨領域的方法論慣例」，
%     絕對不要用它去支撐任何關於 pixel-diff 誤報率的數字。

Accordingly, this study adopts false positive rate on the normal class as
its headline metric, reported alongside anomaly recall to guard against a
degenerate detector that suppresses everything. Following the reporting
convention adopted throughout this project, every metric is reported with
a measure of dispersion and an accompanying interval estimate.
% [VERIFY] 這句話對應你的 Statistic Display Rule。請確認 Chapter 4
%   的實際呈現方式與此處承諾一致（mean ± SD [95% CI] (n = N)），
%   否則審查員會抓到前後不一致。


% ---------------------------------------------------------------------
% 2.8 Gap
% ---------------------------------------------------------------------
\section{Research Gap}
\label{sec:gap}

The preceding sections permit the gap to be stated precisely. Existing
VRT research operates on static screenshot pairs and outputs a binary
flag, leaving temporal defects uncovered and the reviewer's interpretive
burden intact. Temporal defect detection with VLMs has been benchmarked in
the gameplay domain, but as a demonstration of failure rather than as a
proposed remedy, and its conclusions do not transfer directly to web UIs,
whose defects are structured, whose backgrounds are static and whose DOM
provides a parallel non-visual description unavailable in rendered game
footage. The most closely related recent work, WUICC-bench, addresses the
web UI domain and demonstrates that semantic models suppress visual noise
far more selectively than pixel comparison, but operates on static image
pairs containing exactly one atomic change, relies on purely visual
evidence, and requires task-specific model training.

Table~\ref{tab:gap} positions the present study against this prior work.

\begin{table}[htbp]
\centering
\caption{Positioning of the present study relative to the most closely
related prior work. The present study is distinguished by its temporal
input, its use of structural evidence alongside visual evidence, and its
reliance on zero-shot inference rather than task-specific training.}
\label{tab:gap}
\begin{tabular}{@{}llll@{}}
\toprule
 & Conventional VRT & WUICC-bench & Present study \\
\midrule
Input          & Screenshot pair   & Screenshot pair   & Interaction video \\
Defect class   & Static            & Static, atomic    & Temporal \\
Evidence       & Pixels            & Pixels (learned)  & Pixels + DOM \\
Output         & Binary flag       & Caption           & Label + type + rationale \\
Model          & None              & Trained IDC       & Zero-shot VLM \\
Cost concern   & Not applicable    & Not addressed     & Explicit objective \\
\bottomrule
\end{tabular}
\end{table}
% [VERIFY] 這張表的最後一欄請對照你 Chapter 3 的實際設計逐格確認，
%   特別是 "Zero-shot VLM" 與 "Label + type + rationale"
%   （後者對應你的 JSON schema: is_buggy / bug_type / confidence / explanation）。


% ---------------------------------------------------------------------
% 2.9 Summary
% ---------------------------------------------------------------------
\section{Summary}
\label{sec:lit-summary}

This chapter has established three things. First, the pixel-level oracle
on which mainstream visual regression testing depends produces
false-positive rates that recent benchmarking places above
\SI{93}{\percent} on non-meaningful changes, and this deficiency is
representational rather than a matter of threshold tuning. Second, while
vision-language models offer the semantic discrimination that the pixel
oracle lacks, applying them naively reproduces the problem in a different
form, with off-the-shelf models oscillating between excessive sensitivity
and excessive conservatism. Third, no existing work addresses the
temporal class of web UI defects with an approach that combines selective
visual input construction with structural evidence. The following chapter
sets out a methodology designed to address this gap.